\chapter{Revisioni}
\section{Prima revisione}
La prima revisione, svolta in data 09/12/2024, ha evidenziato alcune piccole incorrettezze per quanto riguarda:
\begin{itemize}
    \item I titoli dei task, ambigui in quanto leggendoli si aveva l’impressione di star leggendo dei need
    \item Questionari, la domanda 3 introduttiva ( ora non più presente in quanto eliminata), era ridondante sul quesito della frequenza dell’utente nelle infrastrutture verdi, questo perché già viene chiesto nella domanda numero 9, e i due dati non combaciavano.
    \item Storyboards, nello storyboards numero 3 c’era la presenza di una doppia scena, il che poteva creare problemi.
    \item Un’altra correzione di poco rilievo, in quanto sintattica più che altro, è stata fatta sul modo in cui erano stati numerati i task.
\end{itemize}

\section{Seconda revisione}
La seconda revisione, svolta in data 24/01/2025, ha evidenziato errori sulla prototipazione dell’applicazione.
I problemi riscontrati riguardavano sia i prototipi cartacei che quelli creati in digitale; nello specifico, per quanto riguarda i prototipi cartacei ci è stato evidenziato di come i testi di alcune informazioni erano troppo generici, come il nome e la descrizione degli eventi e delle attività disponibili.
Un’altro problema venuto fuori era la mancanza di una pagina, dopo il compimento del task numero 3, che facesse capire all’utente di aver inserito i filtri e di star visualizzando una schermata con gli eventi filtrati.
Il prototipo digitale numero 1 ha riscontrato gli stessi errori appena citati.